\begin{center}
\textsc{Gauss} to \textsc{Gerling}.\quad Göttingen, 7~November 1822.
\end{center}

\par\noindent\hspace*{\parindent}\mbox{.\ .\ .\ .\ .\ .}\quad At very shallow incidence, where it eventually becomes impossible to direct the beam, I have always had double reflection used with the \textls[120]{most splendid} success, feeding the heliotrope with a hand mirror set up in a suitable place on the ground instead of with the sun itself.
I recommend that you try this technique yourself.\quad\mbox{.\ .\ .\ .\ .\ .}

\par\noindent\hspace*{\parindent}Surveying in mountainous country is a pleasure, and the greater the distances, the greater the pleasure. With heliotropes there is no limit at all to those distances.
My \textls[120]{large} mirror, 1 square foot, had even penetrated the dense smoke from a moorland fire, a pall covering several square miles, from Lichtenberg to Falkenberg, nearly 12 miles. The smoke, to be sure, was too dense for the tiny heliotrope mirrors.\quad\mbox{.\ .\ .\ .\ .\ .}
